\documentclass[11pt]{article}
\usepackage[margin=1in]{geometry}
\usepackage{amsmath,amssymb,amsthm,booktabs,array,microtype}
\usepackage[T1]{fontenc}
\usepackage{lmodern}
\usepackage{xcolor}
\usepackage[hidelinks]{hyperref}
\hypersetup{
  pdftitle={An Exact Audit of an Announced Three-Dimensional Keller Map},
  pdfauthor={Independent audit contributors},
  pdfsubject={Exact verification and consequences of an announced Keller map},
  pdfkeywords={Jacobian Conjecture, Keller map, polynomial map, Weyl algebra,
  Poisson algebra, reproducible mathematics}
}
\newtheorem{theorem}{Theorem}[section]
\newtheorem{proposition}[theorem]{Proposition}
\newtheorem{corollary}[theorem]{Corollary}
\newtheorem{lemma}[theorem]{Lemma}
\theoremstyle{remark}
\newtheorem{remark}[theorem]{Remark}
\newcommand{\A}{\mathbb A}
\newcommand{\C}{\mathbb C}
\newcommand{\R}{\mathbb R}
\newcommand{\Q}{\mathbb Q}
\newcommand{\Pic}{\operatorname{Pic}}
\newcommand{\Gal}{\operatorname{Gal}}
\newcommand{\Disc}{\operatorname{Disc}}
\title{An Exact Audit of an Announced\\Three-Dimensional Keller Map}
\author{Independent audit contributors\\
\small AI-assisted, user-directed working note; maintainer attribution pending}
\date{Revised 21 July 2026}

\begin{document}
\maketitle

\begin{abstract}
We independently audit the polynomial map $F=(P,Q,R):\A^3\to\A^3$
announced by Levent Alpoge on 20 July 2026. Exact calculation gives
$\det DF=-2$, while three distinct rational points map to $(-1/4,0,0)$.
Thus the map is a counterexample to the Jacobian Conjecture in dimension
three. We derive two inverse cubics, prove generic degree three and full
$S_3$ monodromy, classify the complex fibers, determine the image and
nonproperness set, describe a finite smooth completion, and construct direct
nonautomorphic endomorphisms of the third Weyl and Poisson algebras. We also
exhibit a weighted scaling symmetry and collision family, a commuting
divergence-free polynomial frame with no complete constant direction, a
nonproper smooth fibration with jumping component count, an equivariant
quotient square law, exact finite-field fiber statistics, and exact real volume
formulas. We separate proved consequences from open directions and
priority-sensitive claims. Exact symbolic scripts accompany this note. This is
an independent, non-peer-reviewed audit and makes no claim of discovery or
priority.
\end{abstract}

\section{Source, status, and the exact counterexample}

The primary public source located by our search is Alpoge's announcement
\cite{alpoge}. The post credits Akhil for the question and Fable for work on
the example. Several structural observations also appeared in same-day public
notes \cite{gallagher,mo,zhang}; we therefore make no priority claim. As of the
date above we did not locate a peer-reviewed paper or an arXiv manuscript by
the announcer.

Define
\begin{align*}
P&=(1+xy)^3z+y^2(1+xy)(4+3xy),\\
Q&=y+3x(1+xy)^2z+3xy^2(4+3xy),\\
R&=2x-3x^2y-x^3z.
\end{align*}

\begin{theorem}\label{thm:core}
The map $F=(P,Q,R)$ satisfies $\det DF=-2$ and
\[
F(0,0,-1/4)=F(1,-3/2,13/2)=F(-1,3/2,13/2)=(-1/4,0,0).
\]
Its component degrees are $7,6,4$.
\end{theorem}

\begin{proof}
Differentiate and expand the determinant, then substitute the three points.
These are identities in $\mathbb Z[x,y,z,1/2]$. The accompanying verifier
performs the calculation using exact symbolic arithmetic.
\end{proof}

\begin{corollary}
The Jacobian Conjecture is false in dimension three and in every dimension
$n\geq3$. Replacing $P$ by $-P/2$ normalizes the determinant to one.
The result holds over every characteristic-zero field. It does not decide the
two-dimensional case.
\end{corollary}

\subsection{Weighted equivariance and a collision family}

For $\lambda\in\C^*$, define
\[
\sigma_\lambda(x,y,z)=\left(\lambda x,\frac y\lambda,
\frac z{\lambda^2}\right),\qquad
\tau_\lambda(p,q,r)=\left(\frac p{\lambda^2},\frac q\lambda,
\lambda r\right).
\]
Exact substitution gives $F\circ\sigma_\lambda=\tau_\lambda\circ F$. Hence
for every nonzero $\lambda$,
\begin{align*}
F\left(0,0,-\frac1{4\lambda^2}\right)
&=F\left(\lambda,-\frac3{2\lambda},\frac{13}{2\lambda^2}\right)\\
&=F\left(-\lambda,\frac3{2\lambda},\frac{13}{2\lambda^2}\right)
=\left(-\frac1{4\lambda^2},0,0\right).
\end{align*}
The discriminant below is equivariant:
$\Delta(\tau_\lambda(p,q,r))=\lambda^{-2}\Delta(p,q,r)$. Since
$\Delta(-1/4,0,0)=4$, the announced collision is a regular generic
three-sheet fiber, not an exceptional discriminant fiber. The symmetry
$F(-x,-y,z)=(P,-Q,-R)$ is the case $\lambda=-1$.

\subsection{The two-dimensional weighted quotient}

The source and target invariant rings are
\[
\C[x,y,z]^{\mathbb G_m}=\C[u,v],\quad u=xy,\ v=x^2z,
\qquad
\C[p,q,r]^{\mathbb G_m}=\C[\alpha,\beta],\quad
\alpha=pr^2,\ \beta=qr.
\]
Every invariant monomial factors through these generators. Equivariance thus
induces $\bar F:\A^2_{u,v}\to\A^2_{\alpha,\beta}$. Put $h=2-3u-v$.
Exact elimination gives
\begin{align*}
\alpha&=(1+u)h^2(3u^3+u^2v+4u^2+2uv+v),\\
\beta&=h(9u^3+3u^2v+12u^2+6uv+u+3v),
\end{align*}
and
\[
\det D\bar F=-2h^2.
\]
The distinct points $(0,0)$ and $(-3/2,13/2)$ both map to the quotient origin;
indeed the whole critical line $h=0$ maps there. Thus the quotient retains a
shadow of noninjectivity but introduces ramification. It is not a plane Keller
counterexample, and the formula pinpoints why this natural reduction fails.

\begin{theorem}[Equivariant quotient square law]\label{thm:squarelaw}
Let $G=(A,B,C):\A^3\to\A^3$ be equivariant for source weights
$(1,-1,-2)$ and target weights $(-2,-1,1)$, and suppose $\det DG=c\in\C^*$.
With $u=xy$, $v=x^2z$, $\alpha=AC^2$, and $\beta=BC$, write
$C=xh(u,v)$. Then the induced quotient map satisfies
\[
\det D\bar G=c\,h(u,v)^2.
\]
Thus a nonconstant $C/x$ forces a critical divisor in the quotient.
\end{theorem}
\begin{proof}
For the infinitesimal orbit fields $\xi=(x,-y,-2z)$ and
$\eta=(-2A,-B,C)$, direct exterior algebra gives
\[
du\wedge dv=x^2\iota_\xi(dx\wedge dy\wedge dz),\qquad
d\alpha\wedge d\beta=C^2\iota_\eta(dA\wedge dB\wedge dC).
\]
Equivariance gives $DG(\xi)=\eta\circ G$, and the constant Jacobian pulls the
target volume form back to $c$ times the source volume form. Pullback of the
second identity and comparison with the first gives
$\det D\bar G=c(C/x)^2=ch^2$.
\end{proof}

This is a theorem for the displayed weight pattern, not for arbitrary
quotients or every possible descent to dimension two.

\section{Inverse cubics and generic degree}

Write target coordinates as $(p,q,r)$. In the function field, set
$t=y+1/x$. Direct simplification gives
\[
q=4t+\frac2x-3rt^2,\qquad 2p=rt^3-2t^2+qt.
\]
Thus $t$ is a root of
\begin{equation}\label{eq:tcubic}
C_{p,q,r}(T)=rT^3-2T^2+qT-2p.
\end{equation}
If $\rho=C'(t)$, then
\begin{equation}\label{eq:trec}
\rho=\frac2x,\quad x=\frac2\rho,\quad y=t-\frac\rho2,
\quad z=\frac54\rho^2-\frac32t\rho-\frac r8\rho^3.
\end{equation}
Consequently $\C(x,y,z)=\C(p,q,r)(t)$ and the field degree is at most
three. The three simple points in Theorem~\ref{thm:core} have disjoint analytic
neighborhoods mapped biholomorphically onto a common neighborhood. Hence the
generic degree is at least three, and therefore exactly three. In particular,
\eqref{eq:tcubic} is irreducible over $\C(p,q,r)$.

A projectively complete form is obtained from $s=x/(1+xy)$:
\begin{equation}\label{eq:scubic}
H_{p,q,r}(s)=2ps^3-qs^2+2s-r=0,
\end{equation}
with homogenization
\[
\Phi_{p,q,r}(S,T)=2pS^3-qS^2T+2ST^2-rT^3.
\]
For a finite simple root $s$, put $d=1-sq+3ps^2=H'(s)/2$. Then
\[
x=\frac{s}{d},\qquad y=q-3ps,
\qquad z=\frac{2x-3x^2y-r}{x^3}\quad(s\neq0).
\]
For $s=0$, necessarily $r=0$, and the regular formula is
$(x,y,z)=(0,q,p-4q^2)$. For the projective root $[S:T]=[1:0]$,
simplicity forces $p=0$ and $q\neq0$, and
\[
x=\frac2q,\qquad y=-\frac q2,\qquad z=\frac{5x-r}{x^3}.
\]
Direct substitution checks both exceptional formulas. Conversely an affine
point with $1+xy\neq0$ gives $H'(s)=2/(1+xy)\neq0$, while a point with
$1+xy=0$ gives a simple root at infinity. Thus simple projective roots are
in bijection with affine preimages.

\section{Discriminant, monodromy, fibers, and image}

The discriminant of \eqref{eq:tcubic} is $-4\mathcal Q$, where
\[
\mathcal Q=27p^2r^2-18pqr+16p+q^3r-q^2,
\qquad
\Disc_p(\mathcal Q)=-4(3qr-4)^3.
\]
Gauss's lemma shows that $\mathcal Q$ is irreducible and not a square in
$\C(p,q,r)$. Since the cubic is irreducible,
\[
\Gal(C/\C(p,q,r))\cong S_3.
\]
The degree-three cover therefore has full $S_3$ monodromy. Its rational deck
group is trivial because for the nonnormal cubic subextension
$L/K$, $\operatorname{Aut}_K(L)=N_{S_3}(S_2)/S_2=1$.

For \eqref{eq:scubic}, write
\[
\Delta=q^2-16p-q^3r+18pqr-27p^2r^2.
\]
Its discriminant is $4\Delta$. A cubic has no simple projective root precisely
when it has a triple root. Coefficient comparison gives the triple-root curve
\[
\Gamma=\left\{\left(\frac1{3u^2},\frac2u,\frac{2u}3\right):u\in\C^*\right\}
=V(12p-q^2,3qr-4).
\]

\begin{theorem}\label{thm:fibers}
For $(p,q,r)\in\C^3$,
\[
\#F^{-1}(p,q,r)=
\begin{cases}
3,&\Delta\neq0,\\
1,&\Delta=0\text{ and }(p,q,r)\notin\Gamma,\\
0,&(p,q,r)\in\Gamma.
\end{cases}
\]
Hence $F(\C^3)=\C^3\setminus\Gamma$.
\end{theorem}

\begin{proof}
Off $V(\Delta)$ the projective cubic has three simple roots. On the
discriminant but off $\Gamma$, it has a double root and a simple root. The
simple root reconstructs one affine point; the double root cannot, because
$F$ is etale and has no ramified affine point. On $\Gamma$ the unique root is
triple, so no affine preimage exists.
\end{proof}

Generic degree three is minimal for any Keller counterexample. Degree one is
the birational case; a separable degree-two extension is Galois, and the
Galois case of the Jacobian Conjecture is known
\cite{campbell,razar,wright1981}.

\section{Infinity and the smooth finite completion}

The explicit path
\[
(x,y,z)=\left(\lambda,-\frac1\lambda,\frac5{\lambda^2}\right),
\qquad F(x,y,z)=\left(0,\frac2\lambda,0\right)
\]
proves nonproperness.

Define
\[
I=\{((p,q,r),[S:T])\in\A^3\times\mathbb P^1:\Phi_{p,q,r}(S,T)=0\}.
\]
The projection $\pi:I\to\A^3$ is finite of degree three. The partial
derivatives with respect to $p,q,r$ are $2S^3,-S^2T,-T^3$, which never vanish
simultaneously, so $I$ is smooth. The map
\[
(x,y,z)\longmapsto(F(x,y,z),[x:1+xy])
\]
identifies $\A^3$ with $I\setminus D$, where $D$ is the multiple-root divisor.

On the chart $S=1$, solving $\Phi=\partial\Phi/\partial u=0$ gives
\[
p=u^2-ru^3,\qquad q=4u-3ru^2.
\]
Thus $D\cong\A^2_{u,r}$ and its finite birational map to
$\Sigma=V(\Delta)$ is the normalization. The triple-root locus $3ru=2$ maps
to $\Gamma$. Standard cubic-discriminant geometry gives an ordinary transverse
cusp along $\Gamma$.

\begin{proposition}
The full nonproperness set is $S_F=\Sigma$.
\end{proposition}
\begin{proof}
Approaching $D$ from $I\setminus D$ produces escaping source sequences with
limits in $\pi(D)=\Sigma$. Conversely, properness of the finite map
$I\to\A^3$ excludes any other finite asymptotic value.
\end{proof}

Projection $I\to\mathbb P^1$ is a nontrivial affine-plane torsor under
\[
E=\ker\left(\mathcal O^{\oplus3}\xrightarrow{(2S^3,-S^2T,-T^3)}
\mathcal O(3)\right)\cong\mathcal O(-1)\oplus\mathcal O(-2).
\]
The displayed splitting follows from syzygies $(T,2S,0)$ and
$(0,T^2,-S^2)$. There is no global section because constants $p,q,r$ cannot
produce the missing $ST^2$ term. Hence $\Pic(I)\cong\mathbb Z$. Adjunction and
the ramification formula give $K_I\cong\mathcal O(1)\cong\mathcal O(D)$, so
$D$ generates the divisor class group. Thus $I$ is smooth affine but not
factorial, whereas $I\setminus D\cong\A^3$ is factorial.

The existence of some finite completion is classical: Zariski's Main Theorem
factors every separated quasi-finite morphism as an open immersion followed by
a finite morphism \cite{stacks}. The special content here is the explicit boundary,
ramification, normalization, and torsor geometry, not the abstract existence
of a completion.

\section{Direct Dixmier and Poisson counterexamples}

Let $J=DF$ and $B=J^{-1}=-\tfrac12\operatorname{adj}(J)$. Define
\[
\delta_i=\sum_jB_{ji}\frac{\partial}{\partial x_j}.
\]
Then $\delta_i(F_k)=\delta_{ik}$. The commutator
$[\delta_i,\delta_j]$ annihilates $\C(F_1,F_2,F_3)$ and hence its finite
separable extension $\C(x_1,x_2,x_3)$, so the vector fields commute.

For the third Weyl algebra, the assignments
\[
X_i\mapsto F_i(X),\qquad D_i\mapsto\sum_jB_{ji}(X)D_j
\]
preserve all Weyl relations. They do not define an automorphism. Otherwise an
automorphism would take the centralizer $\C[X]$ of $\C[X]$ onto the centralizer
of $\C[F]$. Since every polynomial in $X$ already centralizes $\C[F]$, this
would force $\C[X]\subseteq\C[F]$, and hence equality, contradicting
Theorem~\ref{thm:core}. Thus the Dixmier Conjecture fails in rank three.

On the canonical polynomial Poisson algebra, the assignments
\[
x_i\mapsto F_i(x),\qquad \xi_i\mapsto\sum_jB_{ji}(x)\xi_j
\]
preserve brackets. The same centralizer argument proves noninvertibility, so
the Poisson Conjecture fails in rank three. These implications are consistent
with the classical equivalence framework \cite{ae2007}.

\subsection{A commuting volume-preserving frame and its incompleteness}

The columns of $B$ satisfy the additional exact identities
\[
[\delta_i,\delta_j]=0,\qquad
\operatorname{div}(\delta_i)=0,\qquad
\det(B)=-\frac12.
\]
The divergence equations are the Piola identity for the adjugate of a gradient
matrix and are checked directly by the verifier. Thus the $\delta_i$ form an
everywhere independent commuting, standard-volume-preserving polynomial frame
dual to $dP,dQ,dR$.

\begin{proposition}[Completeness criterion]
For a complex Keller map $G:\A^n\to\A^n$, let $\partial_i$ be the commuting
polynomial fields dual to $dG_i$. Then $G$ is a polynomial automorphism if and
only if all $\partial_i$ are complete holomorphic vector fields. In
particular, local nilpotence of all the $\partial_i$ implies invertibility.
\end{proposition}
\begin{proof}
For an automorphism the fields are pullbacks of constant coordinate fields.
Conversely, their complete commuting flows give a $\C^n$-action satisfying
\[
G(\exp(s_1\partial_1)\cdots\exp(s_n\partial_n)x)=G(x)+s.
\]
The frame condition makes every orbit open, so connectedness gives one orbit;
the formula makes the action free and $G$ bijective. Ax--Grothendieck makes
the inverse polynomial.
\end{proof}

\begin{theorem}[Total directional incompleteness]
Every nonzero constant linear combination
$\delta_a=a_1\delta_1+a_2\delta_2+a_3\delta_3$ is holomorphically incomplete
and is not locally nilpotent.
\end{theorem}
\begin{proof}
A local flow satisfies $F(\phi_t(x))=F(x)+ta$. The omitted curve
\[
\Gamma=\left\{\left(q^2/12,q,4/(3q)\right):q\in\C^*\right\}
\]
contains no affine line. Choose $\gamma\in\Gamma$ and $t_0\ne0$ such that
$y=\gamma-t_0a\notin\Gamma$. Since $F(\A^3)=\A^3\setminus\Gamma$, choose
$x$ with $F(x)=y$. A complete flow would satisfy
$F(\phi_{t_0}(x))=\gamma$, a contradiction.
\end{proof}

An explicit integral curve of $\delta_2$ is
\[
\gamma(t)=\left(2/t,-t/2,5t^2/4\right),\qquad
F(\gamma(t))=(0,t,0),\qquad \gamma'(t)=\delta_2(\gamma(t)).
\]
It escapes at the finite time $t=0$. Thus the exact coframe defines a flat
algebraic affine structure whose developing map is $F$, but every nonzero
constant direction in its dual frame is incomplete.

\section{Real fibers, a plane slice, and research directions}

For real targets, the sign of $\Delta$ gives three real preimages when
$\Delta>0$, one when $\Delta<0$, one on $\Sigma\setminus\Gamma$, and none on
$\Gamma(\R)$. Hence $F(\R^3)=\R^3\setminus\Gamma(\R)$.

If $E\subset\{\Delta>0\}$ is relatively compact and measurable, the three
real inverse branches and $|\det DF|=2$ give
\[
\operatorname{vol}(F^{-1}(E))=\frac32\operatorname{vol}(E).
\]
For relatively compact measurable $E\subset\{\Delta<0\}$, one real branch
gives
\[
\operatorname{vol}(F^{-1}(E))=\frac12\operatorname{vol}(E).
\]
Every local orientation contribution is $-1$, while the signed sum over a
regular real fiber changes from $-3$ to $-1$ when two branches escape at the
discriminant. This is not a global topological degree, since $F$ is not
proper.

Fix $P=p\neq0$. The level surface is
$X=\A^2\setminus V(1+xy)$, and the remaining coordinates are
\[
Q_p=\frac{3px+xy^2+y}{1+xy},\qquad
R_p=\frac{x(-px^2+x^2y^2+3xy+2)}{(1+xy)^3},
\]
with relative Jacobian $-2/(1+xy)^3$. This is an etale map from a punctured
affine plane, not a counterexample on $\A^2$; the nonconstant unit $1+xy$
shows $X\not\cong\A^2$.

There is a second nonproperness diagnostic. Put
\[
A=1+xy,\qquad B=A^2z+y^2(4+3xy),\qquad P=AB.
\]
Since $DF$ is invertible, $dP$ never vanishes and
$P:\A^3\to\A^1$ is smooth. Nevertheless,
\[
P^{-1}(p)\cong X\quad(p\neq0),\qquad
P^{-1}(0)=\{B=0\}\sqcup\{A=0\}
\cong X\sqcup(\mathbb G_m\times\A^1).
\]
The zero-fiber components are disjoint because $A=0$ forces $B=y^2\neq0$.
Thus a smooth polynomial family has connected nonzero fibers but a
disconnected zero fiber; the extra component enters from infinity.

\subsection{Exact finite-field fiber statistics}

Over every finite field $\mathbb F_q$ of odd characteristic, let $N_j(q)$ be
the number of targets having exactly $j$ rational preimages. Only
$j=0,1,3$ occur.

\begin{theorem}\label{thm:finitefields}
If $\operatorname{char}\mathbb F_q\ne2,3$, then
\[
N_3=\frac{(q-1)(q^2+2)}6,\quad
N_1=\frac{q^3+q^2-2q+2}2,\quad
N_0=\frac{(q-1)(q^2+2)}3.
\]
In characteristic three,
\[
N_3=\frac{q^2(q-1)}6,\quad
N_1=\frac{q^2(q+1)}2,\quad
N_0=\frac{q^2(q-1)}3.
\]
The limiting proportions for three, one, and zero preimages are
$1/6,1/2,1/3$, the proportions of the three cycle types in $S_3$.
\end{theorem}
\begin{proof}
A target determines the binary cubic
$2pS^3-qS^2T+2ST^2-rT^3$. Three rational preimages correspond to a
three-element root subset of $\mathbb P^1(\mathbb F_q)$ whose $ST^2$
coefficient is nonzero, since that coefficient must be scaled to $2$.
Of the $\binom{q+1}{3}$ root triples, $(q-1)/2$ bad triples contain infinity
and finite roots $a,-a$. For three finite roots, the bad equation is
$ab+ac+bc=0$; inversion turns it into three distinct nonzero elements with
sum zero. Ordered-pair counting gives $(q-1)(q-5)/6$ bad finite triples when
the characteristic is not three, and $(q-1)(q-3)/6$ in characteristic three.
Subtraction gives $N_3$. Since both source and target have $q^3$ points,
\[
N_0+N_1+N_3=q^3,\qquad N_1+3N_3=q^3,
\]
which gives $N_0=2N_3$ and the stated formulas.
\end{proof}

The immediate research program is now sharply delimited:
\begin{enumerate}
\item Determine whether total-degree $4$, $5$, or $6$ counterexamples exist in
dimension three. Degree at most three is known \cite{wright1993,vistoli}.
\item Produce and optimize explicit cubic-homogeneous and Druzkowski forms via
the genuine stable reductions \cite{bcw}; a naive circuit lift is invalid.
\item Classify the affine-plane torsors and ramification complements producing
$\A^3$.
\item Classify weighted-lift families and their possible monodromy groups. A
same-day note constructs every generic degree at least three \cite{gallagher};
we do not claim priority for that family.
\item Seek explicit small witnesses for related failed conjectures. Published
implications give failure of the Mathieu conjecture for $SU(3)$ and
existential failure in some finite dimension of the Gaussian Moments,
Vanishing, and Image conjectures \cite{mathieu,zwart,gaussian,zhao,image}.
\item Study whether the deleted-hyperbola slice yields an obstruction or a
useful deformation for the still-separate plane problem.
\item Classify the pole orders, escape divisors, and degree growth of the
incomplete inverse-Jacobian flows.
\item Use the weighted $\C^*$-action to form quotient surfaces and classify
equivariant Keller maps.
\item Formalize the determinant, collision, and inverse cubic in a proof
assistant.
\item Extend Theorem~\ref{thm:finitefields} to extension-field zeta functions,
arithmetic monodromy, higher-degree lift families, and bad reduction.
\end{enumerate}

An earlier informal assertion that the map was ``forced in a natural minimal
PDE ansatz'' is deliberately omitted: the class and equivalence relation were
not specified precisely enough for a publication-grade minimality theorem.

\section{Reproducibility and limitations}

The repository verifies the determinant, collision, weighted equivariance,
collision family, weighted quotient, the quotient square-law differential
identities, inverse identities, discriminants, boundary parametrization,
nonproper path, slice Jacobian, first-coordinate factorization, the explicit
escaping integral curve, finite-field regression counts, and commuting
divergence-free inverse-Jacobian fields using exact symbolic arithmetic. These checks
certify polynomial identities, not literature priority or community acceptance.
This note is not peer reviewed, not an arXiv submission, and not a substitute
for independent expert review. Its status discussion is time-stamped and may
be superseded.

\small
\begin{thebibliography}{99}
\bibitem{alpoge} L. Alpoge, announcement on X, 20 July 2026,
\url{https://x.com/__alpoge__/status/2079028340955197566}.
\bibitem{bcw} H. Bass, E. H. Connell, and D. Wright, The Jacobian Conjecture:
Reduction of Degree and Formal Expansion of the Inverse, \emph{Bull. Amer.
Math. Soc.} 7 (1982), 287--330. DOI: 10.1090/S0273-0979-1982-15032-7.
\bibitem{campbell} L. A. Campbell, A Condition for a Polynomial Map to Be
Invertible, \emph{Math. Ann.} 205 (1973), 243--248.
\bibitem{razar} M. Razar, Polynomial Maps with Constant Jacobian, \emph{Israel
J. Math.} 32 (1979), 97--106.
\bibitem{wright1981} D. Wright, On the Jacobian Conjecture, \emph{Illinois J.
Math.} 25 (1981), 423--440.
\bibitem{wright1993} D. Wright, The Jacobian Conjecture: Linear
Triangularization for Cubics in Dimension Three, \emph{Linear Multilinear
Algebra} 34 (1993), 85--97.
\bibitem{vistoli} A. Vistoli, The Jacobian Conjecture in Dimension 3 and Degree
3, \emph{J. Pure Appl. Algebra} 142 (1999), 79--89.
\bibitem{ae2007} P. K. Adjamagbo and A. van den Essen, A Proof of the
Equivalence of the Dixmier, Jacobian and Poisson Conjectures, \emph{Acta Math.
Vietnam.} 32 (2007), 205--214.
\bibitem{zhao} W. Zhao, Hessian Nilpotent Polynomials and the Jacobian
Conjecture, arXiv:math/0409534.
\bibitem{gaussian} H. Derksen, A. van den Essen, and W. Zhao, The Gaussian
Moments Conjecture and the Jacobian Conjecture, \emph{Israel J. Math.} 219
(2017), 917--928.
\bibitem{image} A. van den Essen, D. Wright, and W. Zhao, On the Image
Conjecture, arXiv:1008.3962.
\bibitem{zwart} K. Zwart, Mathieu's Approach to the Jacobian Conjecture,
arXiv:2511.16561.
\bibitem{mathieu} O. Mathieu, Some Conjectures about Invariant Theory and
Their Applications, in \emph{Alg\`ebre non commutative, groupes quantiques et
invariants}, S\'eminaires et Congr\`es 2, Soci\'et\'e Math\'ematique de France,
1997, 263--279.
\bibitem{gallagher} A. Gallagher, Weighted Lifts from the Jacobian
Counterexample, same-day public working note, 2026,
\url{https://github.com/algal/jacobianfun/blob/main/RESEARCH.md}.
\bibitem{mo} Galois Structure of the New Counterexample to the Jacobian
Conjecture, MathOverflow question 513387, 2026,
\url{https://mathoverflow.net/questions/513387/}.
\bibitem{zhang} Z. Zhang, Direct Consequences of the Three-Dimensional
Counterexample to the Jacobian Conjecture, 2026,
\url{https://zzhang-iu.github.io/papers/direct-consequences-jacobian/}.
\bibitem{stacks} The Stacks Project Authors, Zariski's Main Theorem,
Tag 05W7, \url{https://stacks.math.columbia.edu/tag/05W7}.
\end{thebibliography}

\end{document}
